\chapter{Abstract}
\onehalfspacing

The proliferation of digital media on open networks has made the unauthorized copying, redistribution and tampering of multimedia content a pervasive concern. Digital watermarking and steganography are two complementary information-hiding paradigms that address this concern by embedding ownership marks or covert messages directly inside the cover medium, in a manner that is imperceptible to a casual observer. Although a substantial body of work exists on each individual carrier type, most published systems are restricted to a single modality and to a single watermark per cover, which limits their practical use as a general-purpose proof-of-ownership utility.

This project presents \textbf{DWM2 (Digital Water Marking, version 2)}, a unified Java Swing application that supports information hiding across four distinct carrier types: 256 \(\times\) 256 BMP images, WAV audio, plain-text files, and video containers. The image module implements a novel \emph{row-based least-significant-bit (LSB)} scheme in which an MD5 digest of the user-supplied password deterministically selects (i)~the image row that will host the watermark and (ii)~the starting column within that row. Two non-overlapping XOR folds of the digest yield a row index and a column index in the range $[0, 255]$, allowing up to \textbf{256 independent watermarks} to coexist in the same cover image. A single ASCII character is embedded per pixel using a 3-3-2 bit split across the red, green, and blue channels, and a sentinel character `\$' at column~0 of every occupied row prevents accidental collisions. The audio module embeds an XOR-encrypted, timestamp-stamped payload into the LSBs of WAV samples, terminated by a 16-bit end marker. The text module appends an XOR-encrypted block inside an HTML-style comment, and the video module reuses the image algorithm on the first frame of a video that is re-encoded losslessly with FFV1.

The system is implemented in Java SE~17, exposes its functionality through a single graphical user interface, and exhibits fast embedding and extraction (sub-second on a 256\(\times\)256 image). Experimental analysis shows that the embedding distortion is bounded above by \(\pm 7\) per red/green channel and \(\pm 3\) per blue channel, that the password-to-row mapping is uniformly distributed, and that the row-occupation marker reliably blocks overwrites. The DWM2 system therefore demonstrates that simple, transparent and didactically clear primitives can be combined to build a multi-modal information-hiding tool suitable for academic demonstration, ownership proof and lightweight covert communication.

\textbf{Keywords:} Digital watermarking, steganography, least-significant-bit (LSB), MD5 hash, row-based embedding, multi-modal, Java Swing.
